%% ============================================================
%%  Chapter 6 — Conclusion
%% ============================================================
\chapter{Conclusion}
\label{ch:conclusion}

\section{Summary}

This thesis presented HoneyIQ, a simulation framework for training an
adaptive honeypot defender using Deep Q-Networks against a structured,
intent-driven attacker modelled on the Lockheed Martin Cyber Kill Chain.

The key contributions are:

\begin{itemize}
  \item A Markov-chain attacker that produces realistic network traffic
        via UNSW-NB15-inspired feature distributions, with four intent
        profiles encoding qualitatively distinct escalation strategies.
  \item A DQN agent that learns a five-action honeypot policy from a
        24-dimensional state vector encoding attack type, kill-chain stage,
        threat level, escalation rate, and attacker intent.
  \item A domain-informed reward function incorporating proportional
        responses, late kill-chain urgency amplification, and
        honeypot-specific intelligence bonuses for tarpitting.
  \item A composite threat-level metric that integrates attack severity,
        kill-chain progress, escalation rate, and cumulative attack count
        into a single normalised signal.
  \item Empirical evidence that the trained policy achieves a detection
        rate above 0.8 and a false-positive rate below 0.3 on the training
        intent, with moderate transfer to harder intent profiles.
\end{itemize}

\section{Future Work}

Several directions would extend and strengthen this work:

\begin{enumerate}
  \item \textbf{Multi-intent training.}
        Exposing the DQN to a mixture of intent profiles during training
        would improve cross-intent robustness without requiring intent
        inference.

  \item \textbf{Adaptive attacker.}
        Modelling the attacker as a best-responding agent using
        multi-agent RL or game-theoretic formulations would introduce
        true adversarial dynamics.

  \item \textbf{Real-world feature integration.}
        Fitting feature distributions to the actual UNSW-NB15 records,
        or using a generative model (e.g.\ a conditional VAE) trained on
        real packet data, would reduce the simulation-to-reality gap.

  \item \textbf{Classifier-in-the-loop.}
        Replacing ground-truth attack labels in the DQN state with Random
        Forest classifier outputs would test policy robustness under
        observation noise, a prerequisite for real deployment.

  \item \textbf{Uncertainty-aware policy.}
        Replacing the standard DQN with a distributional RL variant
        (e.g.\ C51 or QR-DQN) would allow the agent to reason about
        aleatoric uncertainty in threat levels, potentially improving
        decisions at threat-band boundaries.

  \item \textbf{Evaluation on real honeypot logs.}
        Replaying captured honeypot sessions through the trained policy
        would provide an out-of-distribution test of generalisation to
        genuine attack traffic.
\end{enumerate}

\section{Closing Remarks}

HoneyIQ demonstrates that reinforcement learning is a viable framework for
adaptive honeypot management.
The combination of a structured attacker model, a domain-informed reward
function, and a stable DQN implementation produces a policy that behaves
in alignment with honeypot best practices: engaging low-severity traffic
passively, tarpitting persistent attackers for intelligence, and escalating
to BLOCK or ALERT as campaigns reach critical kill-chain stages.
The simulation infrastructure is fully self-contained, enabling rapid
iteration on both the attacker model and the defender's learning algorithm,
and provides a foundation for more sophisticated adversarial training
scenarios in future research.
